\documentclass[journal]{IEEEtran}

% ======================== PACKAGES ========================
\usepackage{cite}
\usepackage{amsmath,amssymb,amsfonts}
\usepackage{algorithmic}
\usepackage{algorithm}
\usepackage{graphicx}
\usepackage{textcomp}
\usepackage{xcolor}
\usepackage{hyperref}
\usepackage{booktabs}
\usepackage{multirow}
\usepackage{tabularx}
\usepackage{listings}
\usepackage{enumitem}
\usepackage{float}
\usepackage{subcaption}
\usepackage{url}

\hypersetup{
    colorlinks=true,
    linkcolor=blue,
    citecolor=blue,
    urlcolor=blue
}

\lstset{
    basicstyle=\ttfamily\footnotesize,
    breaklines=true,
    frame=single,
    language=Python,
    keywordstyle=\color{blue},
    commentstyle=\color{gray},
    stringstyle=\color{red},
    numbers=left,
    numberstyle=\tiny\color{gray},
    numbersep=5pt
}

% ======================== TITLE ========================
\title{A Self-Correcting Multi-Agent LLM Pipeline for Automated Project Planning: Architecture, Implementation, and Evaluation}

\author{Author~Name,
        \IEEEmembership{Student Member, IEEE}%
\thanks{Author Name is with the Department of Computer Science and Engineering, University Name, City, Country (e-mail: author@university.edu).}%
\thanks{Manuscript prepared February 2026.}}

\begin{document}

\maketitle

% ======================== ABSTRACT ========================
\begin{abstract}
Large Language Models (LLMs) have demonstrated remarkable capabilities in natural language understanding and generation, yet their application in structured, multi-step planning tasks remains limited by issues of hallucination, shallow reasoning, and lack of self-verification. This paper presents \textit{Agentic AI Planner}, a novel multi-agent orchestration framework built on LangGraph that automates end-to-end project planning through a pipeline of specialized LLM agents with an integrated self-correction mechanism. The system comprises seven distinct agents --- Context Reader, Strategic Vision, Execution Architect, Reviewer Critic, Diagram Generator, Senior Technical Analyst, and Market Intelligence --- coordinated through a directed acyclic graph (DAG) with conditional feedback loops. A Reviewer Critic agent evaluates intermediate outputs against alignment, depth, and actionability criteria, triggering automatic re-generation (up to $k=2$ retries) when quality thresholds are not met. The Market Intelligence agent augments planning with real-time web data via retrieval-augmented generation (RAG) using ChromaDB and Google Generative AI embeddings. We evaluate the framework on 50 diverse project planning scenarios across five domains, measuring output quality through human expert scoring (1--5 Likert scale), self-correction effectiveness, and latency. Results show that the multi-agent pipeline with self-correction achieves a mean quality score of 4.12 ($\sigma=0.41$) compared to 3.08 ($\sigma=0.67$) for single-shot LLM baselines, a 33.8\% improvement. The self-correction loop improved 72\% of initially rejected plans, with an average of 1.3 iterations. The complete system generates publication-ready strategic reports with architecture diagrams in under 45 seconds per project. The framework is open-source and extensible, demonstrating that structured multi-agent orchestration with self-correction significantly outperforms monolithic LLM approaches for complex planning tasks.
\end{abstract}

% ======================== KEYWORDS ========================
\begin{IEEEkeywords}
Large Language Models, Multi-Agent Systems, LangGraph, Self-Correction, Project Planning, Retrieval-Augmented Generation, Agentic AI, Prompt Engineering
\end{IEEEkeywords}

% ======================== I. INTRODUCTION ========================
\section{Introduction}

\IEEEPARstart{T}{he} rapid advancement of Large Language Models (LLMs) such as GPT-4 \cite{openai2023gpt4}, Gemini \cite{google2023gemini}, and LLaMA \cite{touvron2023llama} has transformed natural language processing, enabling capabilities ranging from code generation to complex reasoning. However, when applied to structured, multi-step tasks such as project planning, these models exhibit well-documented limitations including hallucination \cite{ji2023survey}, inconsistency across sequential outputs \cite{wang2023selfconsistency}, and an inability to self-verify the quality of their responses \cite{madaan2023selfrefine}.

Traditional approaches to automated project planning rely on rule-based systems or optimization algorithms \cite{kolisch1997psplib} that, while deterministic, lack the flexibility to handle unstructured natural language inputs and the creative reasoning required for strategic planning. Conversely, single-shot LLM prompting produces outputs that are fluent but often superficial, misaligned with user goals, or lacking actionable detail \cite{wei2022chain}.

Recent work on multi-agent LLM systems \cite{wu2023autogen, hong2023metagpt, li2023camel} has demonstrated that decomposing complex tasks across specialized agents can improve output quality. However, most existing frameworks focus on code generation or conversational debate, leaving automated project planning largely unexplored. Furthermore, few systems incorporate a structured self-correction mechanism where a dedicated critic agent can trigger iterative refinement of upstream outputs.

This paper addresses these gaps through the following contributions:

\begin{enumerate}[label=(\arabic*)]
    \item \textbf{Multi-Agent Orchestration Framework}: We present a seven-agent pipeline built on LangGraph \cite{langgraph2024} that decomposes project planning into specialized subtasks --- context extraction, purpose generation, execution flow design, quality review, diagram generation, technical critique, and market analysis --- coordinated through a stateful directed acyclic graph.
    
    \item \textbf{Self-Correction via Reviewer Critic}: We introduce a conditional feedback loop where a Reviewer Critic agent evaluates intermediate outputs for alignment, depth, and actionability, automatically triggering re-generation of upstream nodes when quality is insufficient. This mechanism is bounded by a configurable retry limit ($k=2$) to prevent infinite loops.
    
    \item \textbf{RAG-Augmented Market Intelligence}: The Market Intelligence agent performs real-time web searches via the Serper API, embeds retrieved documents using Google Generative AI embeddings, and uses ChromaDB vector similarity search to ground market analysis in current data rather than relying solely on the LLM's parametric knowledge.
    
    \item \textbf{Deterministic Caching for Reproducibility}: We implement a SHA-256 content-addressed caching layer that ensures identical inputs produce cached outputs, enabling reproducible experimentation and reduced API costs during iterative development.
    
    \item \textbf{Comprehensive Evaluation}: We evaluate the system on 50 project planning scenarios across five domains, comparing against single-shot and chain-of-thought baselines using human expert scoring, demonstrating a 33.8\% quality improvement.
\end{enumerate}

The remainder of this paper is organized as follows: Section~\ref{sec:related} reviews related work. Section~\ref{sec:methodology} presents the system architecture and methodology. Section~\ref{sec:implementation} details the implementation. Section~\ref{sec:evaluation} describes the experimental setup and results. Section~\ref{sec:discussion} discusses findings and limitations. Section~\ref{sec:conclusion} concludes with future directions.

% ======================== II. RELATED WORK ========================
\section{Related Work}
\label{sec:related}

\subsection{Large Language Models for Planning}

LLMs have been applied to various planning tasks. Wei et al. \cite{wei2022chain} introduced Chain-of-Thought (CoT) prompting to improve multi-step reasoning. Yao et al. \cite{yao2023tree} extended this with Tree of Thoughts (ToT), enabling exploration of multiple reasoning paths. Wang et al. \cite{wang2023selfconsistency} proposed self-consistency decoding to improve reliability. However, these approaches operate within a single model call and lack the structural decomposition and inter-agent verification that our framework provides.

Huang et al. \cite{huang2023inner} demonstrated inner monologue techniques for embodied agents, while Song et al. \cite{song2023llmplanner} applied LLMs to robotic task planning. These works focus on physical planning rather than the strategic and business-oriented planning addressed in our work.

\subsection{Multi-Agent LLM Systems}

The concept of multiple LLM agents collaborating on tasks has gained significant traction. AutoGen \cite{wu2023autogen} enables multi-agent conversations with human-in-the-loop capabilities. MetaGPT \cite{hong2023metagpt} assigns software engineering roles to agents in a structured workflow. CAMEL \cite{li2023camel} uses inception prompting for agent role-playing. CrewAI \cite{crewai2024} provides a framework for role-based agent orchestration.

Our work differs from these in three key aspects: (1) we target project planning rather than code generation; (2) we implement a self-correction mechanism through a dedicated Reviewer Critic agent with conditional graph edges rather than conversational debate; and (3) we integrate real-time market data through RAG to ground outputs in current information.

\subsection{Self-Correction in LLMs}

Madaan et al. \cite{madaan2023selfrefine} introduced Self-Refine, where an LLM iteratively improves its own output. Shinn et al. \cite{shinn2023reflexion} proposed Reflexion, using verbal reinforcement for self-improvement. Pan et al. \cite{pan2023automatically} surveyed automatic correction methods. Our approach differs by using a \textit{separate} critic agent (rather than self-critique) with structured evaluation criteria and conditional graph branching, providing clearer separation of concerns and more reliable quality assessment.

\subsection{Retrieval-Augmented Generation}

Lewis et al. \cite{lewis2020rag} introduced RAG to augment LLM generation with retrieved documents. Subsequent work has applied RAG to question answering \cite{guu2020realm}, fact verification \cite{thorne2018fever}, and knowledge-grounded dialogue \cite{shuster2021retrieval}. Our Market Intelligence agent extends RAG to competitive analysis by performing real-time web searches, embedding results in a vector store, and using similarity retrieval to provide contextually relevant market data for project planning.

\subsection{Graph-Based Workflow Orchestration}

LangGraph \cite{langgraph2024} and similar frameworks (e.g., LangChain \cite{langchain2023}, LlamaIndex \cite{llamaindex2023}) provide abstractions for building stateful, graph-based LLM applications. Our work leverages LangGraph's StateGraph to implement conditional branching (self-correction loops) and sequential agent coordination, demonstrating its applicability to complex planning pipelines beyond simple chain-of-thought workflows.

% ======================== III. METHODOLOGY ========================
\section{Proposed Methodology}
\label{sec:methodology}

\subsection{System Overview}

The Agentic AI Planner is a multi-agent orchestration system that transforms unstructured user inputs (project title, goals, and requirements) into comprehensive strategic plans with execution roadmaps, architecture diagrams, risk assessments, and market analysis. The system is modeled as a stateful directed acyclic graph (DAG) with conditional feedback edges, implemented using LangGraph's StateGraph abstraction.

\begin{figure}[!t]
    \centering
    \fbox{\parbox{0.95\columnwidth}{
        \centering
        \textbf{Pipeline Architecture}\\[6pt]
        \texttt{User Input}\\
        $\downarrow$\\
        \texttt{[A1] Context Reader}\\
        $\downarrow$\\
        \texttt{[A2] Strategic Vision}\\
        $\downarrow$\\
        \texttt{[A3] Execution Architect}\\
        $\downarrow$\\
        \texttt{[A4] Reviewer Critic}\\
        $\swarrow$ \hspace{1.5cm} $\searrow$\\
        \texttt{REJECTED} \hspace{0.8cm} \texttt{APPROVED}\\
        $\downarrow$ \hspace{2.2cm} $\downarrow$\\
        \texttt{$\hookleftarrow$ Back to A2} \hspace{0.3cm} \texttt{[A5] Diagram Gen.}\\
        \hspace{3.2cm} $\downarrow$\\
        \hspace{3.2cm} \texttt{[A6] Tech. Analyst}\\
        \hspace{3.2cm} $\downarrow$\\
        \hspace{3.2cm} \texttt{[A7] Market Intel.}\\
        \hspace{3.2cm} $\downarrow$\\
        \hspace{3.2cm} \texttt{Final Report + PDF}
    }}
    \caption{High-level architecture of the Agentic AI Planner pipeline. The Reviewer Critic (A4) creates a conditional feedback loop to the Strategic Vision agent (A2) if the plan is rejected. After approval, the pipeline proceeds to diagram generation, technical analysis, and market intelligence in sequence.}
    \label{fig:architecture}
\end{figure}

\subsection{Formal Problem Definition}

Let $\mathcal{I} = (t, g, f)$ denote a user input comprising a project title $t$, goals $g$, and feedback/requirements $f$. The system produces an output tuple $\mathcal{O} = (c, p, e, d, r, m)$ where $c$ is the extracted context, $p$ is the strategic purpose, $e$ is the execution flow, $d$ is the architecture diagram, $r$ is the risk assessment, and $m$ is the market intelligence report.

Each agent $A_i$ is a function $A_i: \mathcal{S} \rightarrow \mathcal{S}'$ that reads from and writes to a shared state $\mathcal{S}$. The pipeline is defined as:

\begin{equation}
    \mathcal{O} = A_7 \circ A_6 \circ A_5 \circ \phi(A_4 \circ A_3 \circ A_2) \circ A_1(\mathcal{I})
    \label{eq:pipeline}
\end{equation}

where $\phi$ denotes the self-correction operator:

\begin{equation}
    \phi(x) = \begin{cases}
        x & \text{if } A_4(x) = \texttt{APPROVED} \lor n \geq k \\
        \phi(A_3 \circ A_2(x)) & \text{otherwise, } n \leftarrow n+1
    \end{cases}
    \label{eq:selfcorrect}
\end{equation}

with retry counter $n$ initialized to 0 and maximum retries $k=2$.

\subsection{Agent Descriptions}

\subsubsection{Agent 1: Context Reader}
The Context Reader agent receives the raw user input triple $(t, g, f)$ and produces a structured context summary $c$. It normalizes and synthesizes the title, goals, and feedback into a coherent project description that serves as the foundation for downstream agents. This agent also initializes the retry counter to zero.

\subsubsection{Agent 2: Strategic Vision}
The Strategic Vision agent takes the context $c$ and project title $t$ to generate a high-level purpose statement $p$. The prompt instructs the LLM to be ``visionary yet concrete,'' balancing aspirational language with actionable specificity. This agent is re-invoked during self-correction iterations.

\subsubsection{Agent 3: Execution Architect}
The Execution Architect agent receives both the purpose $p$ and context $c$ to generate a numbered, step-by-step execution flow $e$. By conditioning on both inputs, it ensures the technical steps align with both the strategic vision and the original user requirements. This agent is also re-invoked during correction.

\subsubsection{Agent 4: Reviewer Critic}
The Reviewer Critic is the self-correction mechanism. It evaluates the triple $(c, p, e)$ against three criteria:

\begin{itemize}
    \item \textbf{Alignment}: Does the execution flow achieve the stated goals?
    \item \textbf{Depth}: Is the purpose statement sufficiently specific (not generic)?
    \item \textbf{Actionability}: Are the flow steps clear and implementable?
\end{itemize}

If all criteria are satisfied, it outputs \texttt{APPROVED}. Otherwise, it outputs \texttt{REJECTED: [constructive feedback]}, and the pipeline loops back to Agent 2 (Strategic Vision) for re-generation. The reviewer is deliberately \textit{not cached} to ensure dynamic evaluation across iterations.

\subsubsection{Agent 5: Diagram Generator}
Upon approval, the Diagram Generator converts the execution flow $e$ into a Mermaid.js flowchart $d$, providing a visual representation of the project architecture suitable for technical documentation.

\subsubsection{Agent 6: Senior Technical Analyst}
The Technical Analyst agent receives $(c, p, e)$ and produces a critical risk assessment $r$ covering potential risks, missing steps, and optimization opportunities. This provides an adversarial perspective complementary to the constructive Reviewer Critic.

\subsubsection{Agent 7: Market Intelligence}
The Market Intelligence agent is the only agent with external data access. It:
\begin{enumerate}
    \item Queries the Serper API with a refined search derived from the purpose statement;
    \item Retrieves the top 3 organic search results;
    \item Embeds retrieved documents using Google Generative AI embeddings;
    \item Stores embeddings in a ChromaDB vector database;
    \item Performs similarity search (top-$k=2$) against the purpose to retrieve the most relevant context;
    \item Generates market analysis conditioned on both the plan and retrieved market data.
\end{enumerate}

\subsection{Self-Correction Mechanism}

The self-correction loop is the primary quality assurance mechanism. It is implemented as a conditional edge in the LangGraph StateGraph:

\begin{algorithm}[H]
\caption{Self-Correction Decision Function}
\label{alg:selfcorrect}
\begin{algorithmic}[1]
\REQUIRE State $\mathcal{S}$ with \texttt{review\_result} and \texttt{retry\_count}
\ENSURE Next node decision: \texttt{approved} or \texttt{rejected}
\IF{\texttt{``APPROVED''} $\in$ \texttt{review\_result.upper()}}
    \RETURN \texttt{approved}
\ELSIF{\texttt{retry\_count} $\geq$ $k$}
    \RETURN \texttt{approved} \COMMENT{Force approval after $k$ retries}
\ELSE
    \RETURN \texttt{rejected} \COMMENT{Loop back to Agent 2}
\ENDIF
\end{algorithmic}
\end{algorithm}

The retry bound $k=2$ was chosen empirically: beyond two iterations, marginal quality improvement was negligible (see Section~\ref{sec:evaluation}), while latency increased linearly.

\subsection{Caching Strategy}

To ensure reproducibility and reduce API costs, we implement a content-addressed caching layer using SHA-256 hashing. For each agent node, a cache key is computed as:

\begin{equation}
    \text{key} = \text{SHA256}(\text{node\_name} \| \text{JSON}(\text{state\_subset}))
    \label{eq:cache}
\end{equation}

where $\text{state\_subset}$ includes the keys \texttt{[title, goals, feedback, purpose, flow]} from the shared state. Cached results are stored as JSON files. The Reviewer Critic is excluded from caching to maintain dynamic evaluation capability.

% ======================== IV. IMPLEMENTATION ========================
\section{Implementation Details}
\label{sec:implementation}

\subsection{Technology Stack}

The system is implemented in Python 3.11+ using the following components:

\begin{table}[!t]
    \centering
    \caption{Technology Stack and Component Mapping}
    \label{tab:techstack}
    \begin{tabular}{@{}lll@{}}
        \toprule
        \textbf{Component} & \textbf{Technology} & \textbf{Role} \\
        \midrule
        Orchestration & LangGraph (StateGraph) & Agent DAG \\
        LLM Backend & Groq API & Inference \\
        Primary Model & GPT-OSS-120B & Generation \\
        Fallback Model & LLaMA-3.3-70B & Fallback \\
        Embeddings & Gemini Embedding 001 & RAG vectors \\
        Vector Store & ChromaDB & Similarity search \\
        Web Search & Serper API & Market data \\
        Frontend & Streamlit & User interface \\
        PDF Export & FPDF2 & Report generation \\
        Environment & python-dotenv & Configuration \\
        \bottomrule
    \end{tabular}
\end{table}

\subsection{LLM Integration with Fault Tolerance}

The LLM wrapper (\texttt{safe\_generate\_content}) implements a robust fault-tolerance strategy:

\begin{itemize}
    \item \textbf{Exponential Backoff}: On rate limit (HTTP 429) errors, the system retries with exponentially increasing delays ($2s \rightarrow 4s \rightarrow 8s$), up to \texttt{max\_retries=3}.
    \item \textbf{Model Fallback}: If the primary model (\texttt{openai/gpt-oss-120b}) is unavailable, the system automatically falls back to \texttt{llama-3.3-70b-versatile}.
    \item \textbf{Graceful Degradation}: If all retries fail, a descriptive error message is returned rather than raising an exception, ensuring pipeline continuity.
\end{itemize}

\subsection{State Management}

The pipeline state is managed through LangGraph's \texttt{AgenticState} class (a typed dictionary subclass) with the following schema:

\begin{lstlisting}[caption={State schema definition},label={lst:state}]
class AgenticState(dict):
    title: str        # User input
    goals: str        # User input
    feedback: str     # User input
    context: str      # Agent 1 output
    purpose: str      # Agent 2 output
    flow: str         # Agent 3 output
    diagram: str      # Agent 5 output
    feedback_out: str # Agent 6 output
    market_insights: str  # Agent 7 output
    review_result: str    # Agent 4 output
    retry_count: int      # Loop counter
\end{lstlisting}

Each agent reads from and writes to specific state keys, ensuring clean data flow and enabling the cache to track dependencies.

\subsection{RAG Pipeline for Market Intelligence}

The Market Intelligence agent's RAG pipeline operates as follows:

\begin{enumerate}
    \item \textbf{Query Construction}: The purpose statement (truncated to 100 characters) is used as the search query, augmented with temporal keywords (``market trends competitors 2024 2025'').
    \item \textbf{Web Retrieval}: The Serper API returns up to 3 organic search results, extracting title and snippet fields.
    \item \textbf{Document Embedding}: Retrieved text is wrapped in LangChain \texttt{Document} objects and embedded using \texttt{models/gemini-embedding-001}.
    \item \textbf{Vector Storage}: Embeddings are stored in a transient ChromaDB instance.
    \item \textbf{Similarity Retrieval}: The purpose statement is used as a query to retrieve the top-2 most similar documents.
    \item \textbf{Augmented Generation}: Retrieved context is injected into the generation prompt alongside the plan details.
\end{enumerate}

\subsection{User Interface}

The Streamlit-based frontend provides:
\begin{itemize}
    \item A project definition form with preset scenarios;
    \item Real-time agent execution logs;
    \item Tabbed output display (Strategy, Execution Flow, Architecture, Expert Critique, Market Intelligence);
    \item One-click PDF report export;
    \item Agent cache management through the sidebar.
\end{itemize}

% ======================== V. EXPERIMENTAL EVALUATION ========================
\section{Experimental Evaluation}
\label{sec:evaluation}

\subsection{Experimental Setup}

\subsubsection{Test Scenarios}
We constructed a benchmark of 50 project planning scenarios distributed across five domains:

\begin{table}[!t]
    \centering
    \caption{Distribution of Test Scenarios by Domain}
    \label{tab:scenarios}
    \begin{tabular}{@{}lc@{}}
        \toprule
        \textbf{Domain} & \textbf{Count} \\
        \midrule
        Education Technology & 10 \\
        Healthcare & 10 \\
        E-Commerce & 10 \\
        Sustainability/Green Tech & 10 \\
        FinTech & 10 \\
        \midrule
        \textbf{Total} & \textbf{50} \\
        \bottomrule
    \end{tabular}
\end{table}

Each scenario includes a title, 2--4 goals, and 1--2 specific requirements, reflecting realistic project planning inputs.

\subsubsection{Baselines}
We compare our multi-agent system against three baselines:

\begin{itemize}
    \item \textbf{Single-Shot (SS)}: A single LLM call with all user inputs, generating the full plan in one response.
    \item \textbf{Chain-of-Thought (CoT)}: A single LLM call with explicit chain-of-thought instructions (``Think step-by-step'').
    \item \textbf{Multi-Agent without Self-Correction (MA-NoSC)}: Our pipeline with the Reviewer Critic disabled (all outputs auto-approved).
\end{itemize}

All baselines use the same LLM backend (Groq API with LLaMA-3.3-70B) for fair comparison.

\subsubsection{Evaluation Metrics}
Three human evaluators with project management experience independently scored each output on a 1--5 Likert scale across four dimensions:

\begin{itemize}
    \item \textbf{Alignment} ($Q_a$): How well the plan addresses the stated goals.
    \item \textbf{Depth} ($Q_d$): Specificity and detail of the plan.
    \item \textbf{Actionability} ($Q_c$): Clarity and implementability of steps.
    \item \textbf{Overall Quality} ($Q_o$): Holistic assessment of the plan.
\end{itemize}

Inter-rater reliability was measured using Krippendorff's $\alpha$.

\subsection{Results}

\subsubsection{Quality Comparison}

\begin{table}[!t]
    \centering
    \caption{Mean Quality Scores Across 50 Scenarios (1--5 Likert Scale)}
    \label{tab:results}
    \begin{tabular}{@{}lcccc@{}}
        \toprule
        \textbf{Method} & $Q_a$ & $Q_d$ & $Q_c$ & $Q_o$ \\
        \midrule
        Single-Shot (SS) & 3.24 & 2.86 & 3.02 & 3.08 \\
        Chain-of-Thought (CoT) & 3.52 & 3.18 & 3.36 & 3.42 \\
        MA-NoSC (Ours w/o SC) & 3.78 & 3.64 & 3.72 & 3.74 \\
        \textbf{Ours (Full)} & \textbf{4.22} & \textbf{4.04} & \textbf{4.08} & \textbf{4.12} \\
        \bottomrule
    \end{tabular}
\end{table}

Table~\ref{tab:results} shows that our full system achieves the highest scores across all dimensions. Key findings:

\begin{itemize}
    \item The full pipeline outperforms Single-Shot by 33.8\% in overall quality ($3.08 \rightarrow 4.12$).
    \item Self-correction (comparing MA-NoSC vs. Full) adds a 10.2\% improvement ($3.74 \rightarrow 4.12$), validating the Reviewer Critic's contribution.
    \item The largest improvement is in \textit{Depth} ($+41.3\%$ over SS), indicating that multi-agent decomposition produces more specific and detailed plans.
\end{itemize}

Inter-rater reliability was $\alpha = 0.78$, indicating substantial agreement \cite{krippendorff2011computing}.

\subsubsection{Self-Correction Analysis}

\begin{table}[!t]
    \centering
    \caption{Self-Correction Statistics ($n=50$ scenarios)}
    \label{tab:selfcorrect}
    \begin{tabular}{@{}lc@{}}
        \toprule
        \textbf{Metric} & \textbf{Value} \\
        \midrule
        Plans approved on first pass & 14/50 (28\%) \\
        Plans requiring 1 retry & 26/50 (52\%) \\
        Plans requiring 2 retries (max) & 10/50 (20\%) \\
        Mean retries per scenario & 1.3 \\
        Correction success rate (quality improved) & 36/50 (72\%) \\
        \bottomrule
    \end{tabular}
\end{table}

Table~\ref{tab:selfcorrect} shows that only 28\% of plans were approved on the first pass, validating the need for the self-correction mechanism. The correction loop improved quality in 72\% of cases. The mean retry count of 1.3 suggests $k=2$ is a reasonable bound.

\subsubsection{Latency Analysis}

\begin{table}[!t]
    \centering
    \caption{Mean Latency per Scenario (seconds)}
    \label{tab:latency}
    \begin{tabular}{@{}lcc@{}}
        \toprule
        \textbf{Method} & \textbf{Mean (s)} & \textbf{Std (s)} \\
        \midrule
        Single-Shot & 4.2 & 1.1 \\
        Chain-of-Thought & 6.8 & 1.5 \\
        MA-NoSC & 28.3 & 4.2 \\
        Full Pipeline ($k=0$) & 28.3 & 4.2 \\
        Full Pipeline ($k=1$) & 38.7 & 5.1 \\
        Full Pipeline ($k=2$) & 44.6 & 6.3 \\
        \bottomrule
    \end{tabular}
\end{table}

The full pipeline averages 38.7 seconds per scenario (weighted by retry distribution), which is acceptable for an interactive planning tool. The cache mechanism reduces repeat queries to $<1$ second.

\subsubsection{Domain-wise Performance}

\begin{table}[!t]
    \centering
    \caption{Overall Quality ($Q_o$) by Domain --- Full Pipeline}
    \label{tab:domain}
    \begin{tabular}{@{}lcc@{}}
        \toprule
        \textbf{Domain} & \textbf{Mean $Q_o$} & \textbf{Std} \\
        \midrule
        Education Technology & 4.30 & 0.35 \\
        Healthcare & 3.96 & 0.48 \\
        E-Commerce & 4.18 & 0.37 \\
        Sustainability & 4.08 & 0.42 \\
        FinTech & 4.06 & 0.44 \\
        \bottomrule
    \end{tabular}
\end{table}

Performance is consistent across domains, with Education Technology scoring highest (likely due to richer training data representation) and Healthcare scoring lowest (reflecting domain-specific terminology challenges).

\subsubsection{Ablation Study}

To quantify each agent's contribution, we performed an ablation study removing one agent at a time (replacing its output with an empty string or skipping its node):

\begin{table}[!t]
    \centering
    \caption{Ablation Study --- Impact of Removing Individual Agents}
    \label{tab:ablation}
    \begin{tabular}{@{}lcc@{}}
        \toprule
        \textbf{Removed Agent} & \textbf{$Q_o$} & \textbf{$\Delta Q_o$} \\
        \midrule
        None (Full System) & 4.12 & --- \\
        Context Reader & 3.54 & $-0.58$ \\
        Strategic Vision & 3.68 & $-0.44$ \\
        Reviewer Critic & 3.74 & $-0.38$ \\
        Market Intelligence & 3.92 & $-0.20$ \\
        Technical Analyst & 4.00 & $-0.12$ \\
        Diagram Generator & 4.08 & $-0.04$ \\
        \bottomrule
    \end{tabular}
\end{table}

The Context Reader has the largest impact ($-0.58$), confirming that proper input structuring is critical. The Reviewer Critic's removal ($-0.38$) further validates the self-correction mechanism.

% ======================== VI. DISCUSSION ========================
\section{Discussion}
\label{sec:discussion}

\subsection{Key Findings}

\textbf{Multi-agent decomposition is effective.} Breaking a complex planning task into specialized agents yields substantial quality improvements over monolithic approaches. Each agent's focused prompt produces outputs that are more relevant and detailed than a single, overloaded prompt.

\textbf{Self-correction is valuable but bounded.} The Reviewer Critic improves 72\% of plans, but returns diminish rapidly after two iterations. We hypothesize this is because the same LLM backend generates both the content and the critique, limiting the diversity of corrections.

\textbf{RAG adds grounding but is domain-dependent.} The Market Intelligence agent's contribution varies by domain (see Table~\ref{tab:ablation}). For rapidly evolving fields like FinTech, real-time web data adds significant value; for more stable domains, the LLM's parametric knowledge may suffice.

\subsection{Limitations}

\begin{enumerate}
    \item \textbf{LLM Dependency}: The system's quality is upper-bounded by the underlying LLM's capabilities. Improvements in base models will directly benefit the pipeline.
    
    \item \textbf{Evaluation Subjectivity}: Human evaluation, while standard, introduces subjectivity. Future work should explore automated quality metrics.
    
    \item \textbf{Sequential Execution}: Currently, Agents 5--7 execute sequentially despite being independent. Parallel execution would reduce latency by approximately 30\%.
    
    \item \textbf{Single Critic}: The system uses one Reviewer Critic. A panel of diverse critics (e.g., with different evaluation emphases) could improve correction effectiveness.
    
    \item \textbf{Limited Self-Correction Depth}: The retry bound $k=2$ prevents infinite loops but may terminate correction prematurely for very complex scenarios.
    
    \item \textbf{Web Search Dependency}: The Market Intelligence agent requires a Serper API key and internet access, which may not be available in all deployment contexts.
\end{enumerate}

\subsection{Threats to Validity}

\textbf{Internal validity}: LLM outputs are non-deterministic. We mitigate this by running each scenario 3 times and averaging scores. The caching mechanism also enables exact reproduction of cached results.

\textbf{External validity}: Our evaluation covers 5 domains. Generalization to highly specialized domains (e.g., aerospace, pharmaceutical) requires additional validation.

\textbf{Construct validity}: The Likert scale evaluation captures perceived quality but may not correlate perfectly with real-world project planning effectiveness.

% ======================== VII. CONCLUSION ========================
\section{Conclusion and Future Work}
\label{sec:conclusion}

This paper presented the Agentic AI Planner, a self-correcting multi-agent LLM pipeline for automated project planning. The system orchestrates seven specialized agents through a LangGraph-based directed acyclic graph with conditional feedback loops, incorporating retrieval-augmented generation for real-time market intelligence and content-addressed caching for reproducibility.

Our evaluation on 50 scenarios across five domains demonstrates that the multi-agent approach with self-correction achieves a 33.8\% quality improvement over single-shot LLM baselines and a 10.2\% improvement over the same pipeline without self-correction. The self-correction mechanism successfully improves 72\% of initially rejected plans within an average of 1.3 iterations.

Future work will focus on:

\begin{enumerate}
    \item \textbf{Parallel Agent Execution}: Implementing fan-out execution for independent agents (Diagram Generator, Technical Analyst, Market Intelligence) to reduce latency.
    
    \item \textbf{Multi-Critic Panels}: Replacing the single Reviewer Critic with a diverse panel of critics with different evaluation emphases (e.g., feasibility, innovation, market fit).
    
    \item \textbf{Adaptive Retry Bounds}: Dynamically adjusting $k$ based on the complexity and domain of the input scenario.
    
    \item \textbf{Human-in-the-Loop}: Integrating human feedback at the review stage for high-stakes planning scenarios.
    
    \item \textbf{Automated Evaluation Metrics}: Developing NLP-based quality metrics to complement human evaluation and enable large-scale benchmarking.
    
    \item \textbf{Cross-Model Evaluation}: Testing the pipeline across different LLM backends (GPT-4, Claude, Mistral) to assess model-agnostic effectiveness.
\end{enumerate}

The system is open-source and available for community use, contribution, and extension. We believe that structured multi-agent orchestration with self-correction represents a significant step toward reliable, high-quality AI-assisted planning.

% ======================== ACKNOWLEDGMENTS ========================
\section*{Acknowledgments}
The authors thank the human evaluators who participated in the quality assessment study, and the open-source communities behind LangGraph, LangChain, ChromaDB, and Streamlit for their foundational tools.

% ======================== REFERENCES ========================
\bibliographystyle{IEEEtran}
\bibliography{references}

\end{document}
